\documentclass[11pt]{article}
\usepackage[T1]{fontenc}
\usepackage{amsmath,amssymb,amsthm,booktabs}
\usepackage[margin=0.9in]{geometry}
\usepackage[colorlinks=true,linkcolor=blue,citecolor=blue,urlcolor=blue]{hyperref}
\setlength{\emergencystretch}{1em}
\newtheorem{theorem}{Theorem}
\newtheorem{lemma}[theorem]{Lemma}
\newcommand{\F}{\mathbf F}
\newcommand{\A}{A(2)}
\title{A support modification in the Duursma--Mak bound for $A(2)$}
\author{Ihara $A(2)$ Lab\\\small Working draft for mathematical review}
\date{13 September 2026}
\begin{document}
\maketitle
\begin{abstract}
Using the genus-$343$ function field and the class-field-tower criterion of
Duursma and Mak, we obtain
\[
 A(2)\ge \frac{196608}{620201}=0.3170069058256919934\ldots.
\]
The change is confined to the support of a conductor-two modulus: two places
of degree $10$ are replaced by places of degrees $8$ and $12$. The certified
genus denominator decreases by $15/1024$, while the sufficient
Golod--Shafarevich inequality retains positive margin. A small computation in
$\F_{64}$ supplies the needed degree-$12$ place. We give the proof, an
exchange identity explaining the gain, and reproducible exact certificates.
This is a parameter refinement of the existing method. Priority has not been
established by the bounded literature search accompanying this draft.
\end{abstract}

\section{The result and the published construction}
For a prime power $q$, let $N_q(g)$ denote the largest number of rational
points on a smooth projective geometrically irreducible curve of genus $g$
over $\F_q$, and let $A(q)=\limsup_{g\to\infty}N_q(g)/g$.
All function fields below have their indicated full constant field.

The displayed binary conductor estimate of Duursma--Mak\cite{DM} gives the
exact fraction $24576/77527$. We refine its parameters as follows.
\begin{theorem}\label{thm:main}
One has $A(2)\ge196608/620201>24576/77527$.
\end{theorem}
The conclusion does not determine $A(2)$, provide a new recursive equation,
or compute an auxiliary ray class group exactly.

Set
\[
 a=x^3+x+1,\qquad b=x^2+x,\qquad c=x^5+x^2+1.
\]
Let $H/\F_2$ be the function field with $y^2+ay=b$, and let $k/H$ be the
compositum of
\begin{equation}\label{eq:cover}
 w_i^2+w_i=bh_i/c^2,\qquad (h_1,\ldots,h_5)=(1,x,x^2,y,y^2).
\end{equation}
These are normalized function fields; affine singularities are not counted as
additional places. Duursma--Mak's construction has
\begin{equation}\label{eq:geometry}
 [k:H]=32,\quad g(H)=2,\quad g(k)=343,\quad N(k)=192,
 \qquad (B_5,B_6,B_8,B_{10})(k)=(2,16,16,64),
\end{equation}
where $B_d$ counts closed places of degree $d$. In particular, one degree-$8$
place is unused by their published modulus.

For clarity, the geometric ingredients behind\eqref{eq:geometry} are the
following. At each of the two degree-$5$ places of $H$ above $c=0$, the reduced
simple-pole residues of the five functions in\eqref{eq:cover} form a map of
rank five over $\F_2$. Thus every nonzero character has conductor
$2P_5+2P'_5$, of degree $20$; there is no other ramification. The five
Artin--Schreier classes remain independent after extending constants. Hence
\[
 2g(k)-2=32(2g(H)-2)+31\cdot20=684.
\]
The six rational places of $H$ split completely. The accompanying M1
certificates independently reconstruct these facts, the place inventory, and
the zeta function with both binary polynomial arithmetic and SageMath.

\section{A degree-twelve place}
\begin{lemma}\label{lem:place}
The field $k$ has at least sixteen places of degree $12$.
\end{lemma}
\begin{proof}
Let $u$ generate $\F_{64}=\F_2[u]/(u^6+u+1)$. The polynomial is irreducible;
this is also verified by the standalone finite-field checker. Put
\[
 x=u,\qquad y=u^4+u^2.
\]
These values satisfy $y^2+(x^3+x+1)y=x^2+x$. Since $x$ has degree six,
the pair defines a degree-$6$ place $Q$ of $H$. Also $c(x)\ne0$, so $Q$ is
unramified in $k/H$.

Its Frobenius in $\operatorname{Gal}(k/H)=(\F_2)^5$ is the vector
\[
 \left(\operatorname{Tr}_{\F_{64}/\F_2}(bh_i/c^2)\right)_{i=1}^5
 =(0,0,1,1,0).
\]
Only nonvanishing is needed. For $h_3=x^2$, the radicand is $u^5+u^4$.
In polynomial-basis bit notation, its six successive squares are
\[
 48,\ 60,\ 47,\ 38,\ 36,\ 32;
\]
their bitwise sum is $1$. This verifies its trace directly.
The Frobenius has order two, so the unramified place $Q$ has $32/2=16$
places above it, all of degree $6\cdot2=12$.
\end{proof}
The larger certificate lists thirteen such degree-$6$ places of $H$ and
therefore supplies $208$ degree-$12$ places of $k$. Direct orbit enumeration
also gives their coordinates. Lemma\ref{lem:place} is sufficient for the theorem.

\section{Proof of the improved bound}
Take $T$ to be all $192$ rational places of $k$. Choose a set $S$ of distinct
closed places with multiplicities as in the second row of the table, and put
$\mathfrak m=2\sum_{P\in S}P$.
\begin{center}\begin{tabular}{lrrrrrr}
\toprule
Selection & $n_5$ & $n_6$ & $n_8$ & $n_{10}$ & $n_{12}$ & GS margin\\
\midrule
Duursma--Mak &2&16&15&4&0&57\\
Modified support &2&16&16&2&1&41\\
\bottomrule
\end{tabular}
\end{center}
The availability follows from\eqref{eq:geometry} and Lemma\ref{lem:place}.
All places in $S$ are nonrational, hence $S\cap T=\varnothing$.

Let $K$ be the maximal abelian $2$-primary subextension of the $T$-ray class
field with modulus $\mathfrak m$. The $2$-primary convention is essential for
the finite $2$-group argument in the tower criterion; it preserves the
$2$-rank estimate, complete splitting of $T$, and the conductor bounds.
Write $n=[K:k]$, leaving its exact value unspecified.

For the binary conductor-two specialization of Duursma--Mak's Theorem 1.1,
write
\begin{align*}
 A&=\sum_{P\in S}\deg P=266,\\
 C&=\sum_{P\in S}\deg P(\deg P+1)=2260,\\
 D&=1+A-|T|=75.
\end{align*}
We have $0<|T|\le A$, and the sufficient infinitude inequality is
\begin{equation}\label{eq:GS}
 D^2-2C-4A=5625-4520-1064=41>0.
\end{equation}
Indeed, if the resulting unramified split $2$-tower $L/K$ were finite, the
published rank estimates for $G=\operatorname{Gal}(L/k)$ would give
$d_2(G)\ge75$ and $r_2(G)-d_2(G)\le C/2+191=1321$.
On $d\ge75$, the increasing function $d^2/4-d$ is at least $5325/4>1321$,
contradicting the strict Golod--Shafarevich inequality for a finite nontrivial
$2$-group. Thus the tower is infinite.

At a place of degree $f$ in $S$, the inertia image in $\operatorname{Gal}(K/k)$
has order at most $2^f$, the size of $U^1/U^2$. A uniform global character
restricts uniformly to the dual of this inertia image, so it is unramified
there with probability at least $2^{-f}$. There are no conductor-one
$2$-primary characters. Consequently its mean local conductor exponent is at
most $2(1-2^{-f})$. This remains true if the actual inertia is smaller;
no independence between different places is assumed.

By linearity, the mean degree $M$ of the global character conductor satisfies
\begin{align*}
 M&\le2\bigl[10(1-2^{-5})+96(1-2^{-6})+128(1-2^{-8})\\
 &\hspace{37mm}+20(1-2^{-10})+12(1-2^{-12})\bigr]
 =\frac{269993}{512}.
\end{align*}
Conductor--discriminant and Riemann--Hurwitz give
\[
 g(K)-1=n\left(342+\frac M2\right)\le n\frac{620201}{1024}.
\]
There are $192n$ rational places of $K$ above $T$, all split in the infinite
unramified tower. Nonempty rational splitting fixes the constant field
$\F_2$ at every level. Since $g(K)>1$, the genera are unbounded and the
usual unramified-tower ratio yields
\[
 A(2)\ge \frac{192n}{g(K)-1}\ge\frac{196608}{620201}.
\]
Finally,
\[
 \frac{196608}{620201}-\frac{24576}{77527}
 =\frac{368640}{48082322927}>0,
\]
proving Theorem\ref{thm:main}. Notice that $n$ cancels; neither its exact
value nor the invariant factors of the ray group have been assumed.

\section{Why the exchange helps}
Define $w(f)=f(1-2^{-f})$, the conductor-two contribution to the certified
genus denominator. The following elementary identity explains the refinement.
\begin{lemma}
Fix an admissible binary conductor-two support and $T$. Suppose two selected
places of degree $d$ can be replaced by unused distinct places of degrees
$d-h,d+h$, disjoint from $T$, with $h\ge1$ and $d-h\ge2$.
The exchange leaves $A,D$ fixed, raises $C$ by $2h^2$, and lowers the GS
margin by $4h^2$. Its decrease in the genus denominator is
\[
 \Delta(d,h)=\frac{(2^h-1)[d(2^h-1)-h(2^h+1)]}{2^{d+h}}.
\]
If the old margin is at least $4h^2$ and
$d>h(2^h+1)/(2^h-1)$, the new certified ratio is strictly larger.
\end{lemma}
\begin{proof}
The degree sum stays $2d$, while
$(d-h)(d-h+1)+(d+h)(d+h+1)-2d(d+1)=2h^2$.
Expanding $2w(d)-w(d-h)-w(d+h)$ gives the stated $\Delta$.
Its sign and the remaining margin imply the conclusion.
\end{proof}
Here $d=10,h=2$: the margin falls from $57$ to $41$, and
$\Delta=15/1024$. The published denominator $77527/128$ becomes
$620201/1024$. This is an improvement of the conductor upper bound used by
the theorem, not a computation of the actual genus of either ray field.
The identity can prune a restricted search only when the replacement remains
inside that search's allowed degrees and multiplicities.

\section{Computation, scope, and coding interpretation}
The parameter search takes a certified place inventory, not the finite ratio
$N/g$, as input. For fixed $S$ its largest admissible splitting size is
\[
 t_{\max}=\min\{N,A,1+A-\lceil\sqrt{2C+4A}\rceil\}.
\]
A nonpositive value gives no certificate. Any score strictly above an incumbent
$b$ must have conductor cost less than $N/b-(g-1)$, allowing exact monotone
pruning. Two implementations found the same multiplicity vector, without
supplying the improved vector as input. In the $k$, nonrational $S$,
degree-at-most-$12$, depth-two domain there are $11\,778\,195$ possible
multiplicity vectors; $99\,835$ survive the cost pruning. The displayed
multiplicity vector and $t$ uniquely maximize this conservative score in
that domain. This is not uniqueness of a divisor or an optimality claim for
all towers. A separate bounded neighborhood of $14\,112$ depth-$2,4,6$
variants did not improve the score.

The binary elementary AG-code estimate is formally $R+\delta\ge1-1/a$ for
a suitable tower ratio $a$. Here it gives $-423593/196608<0$, hence no new
positive binary coding guarantee. No claim about optimal or industrial codes
is made.

\paragraph{Reproducibility and review status.}
The accompanying repository contains equations, local residue maps, all
selected closed-place coordinates, exact fractions, positive and adversarial
controls, independent SageMath checks, and a full tower proof review.
The checks passed with Python 3.12.3 and SageMath 10.9. The note was prepared
in an AI-assisted laboratory using Astra for integration and adversarial proof
review and Luna for bounded arithmetic, implementation, and bibliography.
These internal checks are not external peer review. The literature audit dated
13 September 2026 found no exact antecedent in its recorded public-source
coverage; the novelty label remains \emph{possibly novel}, with no priority
claim or external submission.

\begin{thebibliography}{9}
\bibitem{DM} I. Duursma and K.-H. Mak,
\emph{On lower bounds for the Ihara constants $A(2)$ and $A(3)$},
arXiv:1102.4127v4 (2012), especially Theorem 1.1 and Proposition 5.5.
Published in \emph{Compositio Mathematica} \textbf{149} (2013), 1108--1128,
\href{https://doi.org/10.1112/S0010437X12000796}{doi:10.1112/S0010437X12000796}.
\end{thebibliography}
\end{document}
